
\section{Abstract computation}\label{sec:comp_sampling}
Probabilistic sampling is a powerful and flexible paradigm of general computation. Sampling methods represent and evaluate probability distributions through the production of stochastic samples samples. Uncertain, high-dimensional, and ambiguous inference problems require efficient and accurate characterisations of target distributions. Sampling methods draw from a target distribution such that the sampled ensemble distribution reflects the target itself. This paradigm has become widely utilised in machine learning and statistics, where generative models learn patterns from large multivariate distributions inferred from observational data \cite{delacour_neural_1997, cranmer_advances_2023}. This framework has gained significant traction as a basic scheme to describe biological computation. The position that the brain itself acts as a Bayesian generative model has become that interprets through continuous sampling mechanisms has become well supported.
\subsection{Active inference}\label{subsec:active_inf}
The dominant framework in modern computational neuroscience is that of the Bayesian brain. Karl Friston, the most cited neuroscientist \cite{Bohannon2016-kp}, developed this framework, and shifted the view of the perceptual loop of the brain. Friston's theory of active inference described perception through inferring and generating top-down sensory predictions of the world, and carrying back the bottom-up errors from that perception. This development brought a strong shift in the literature, ``marking a watershed between 20th-century thinking about the brain as a glorious stimulus–response link and more constructivist 21st century perspectives that emphasized an active sampling of the sensory world'' (\cite{friston_does_2018} pg. 1019). The active inference model of an updating internal mental state filled the explanatory gap in the understanding of the functional role of many key feedback mechanisms in the brain \cite{rao_predictive_1999, koch_predicting_1999}, and went on to yield successful predictions of how the brain optimally deciphers noisy multimodal sensory inputs \cite{kording_bayesian_2004}, as well as provide rigorous frameworks to base probabilistic computation in the brain \cite{lange_bayesian_2020}. There has recently been opposition to the general framework of the Bayesian brain and active inference \cite{mangalam_myth_2025}. Much of the arguments against such a framework arise from the philosophy of science, where active inference in perception may be supported post-hoc through ``just so'' definitions of priors \cite{mangalam_myth_2025, bowers_bayesian_2012}. It is therefore important to understand the Bayesian brain hypothesis as a normative framework from which to build falsifiable, understandable, and biologically plausible theories of neural computation.

\subsubsection*{Visual search}
This thesis presents visual search tasks as such a paradigm where a Bayesian model is well motivated. The visual system is evolved to actively and efficiently process noisy information and evidence from the external world, and plan future search paths to gather further noisy evidence \cite{Renninger2005-ee, renninger_where_2007, najemnik_optimal_2005, bujia_modeling_2022}. In particular, search tasks require sophisticated integration of the cost and expected information gain of decisions, which invites fruitful decision analysis \cite{Wedel_2023}. Empirical experiments show that eye movements show good agreement with greedy local information maximising fixations, given well-motivated visual sensitivity based on human physiology \cite{Radulescu_2026, najemnik_optimal_2005, renninger_where_2007}.

\subsection{Neural sampling}\label{subsec:neur_sampling}
Another useful paradigm for the basic functional role of neural activity is the neural sampling hypothesis. The framework of neural sampling posits that the variability of cortical neural responses is an important functional feature of biological computation. Otherwise regarded as noise, stochastic neural activity patterns are thus presented as samples drawn from a posterior distribution over possible sensory interpretations \cite{orban_neural_2016, NIPS2002_a486cd07}. The neural sampling hypothesis has been shown to perform and predict network-level empirical patterns of neural activity in the V1, and recently by Terada and Toyoizumi in biologically realistic cortical networks, better than alternative theories \cite{NIPS2002_a486cd07, ujfalussy_sampling_2022, terada_chaotic_2024, orban_neural_2016}. Perceptual uncertainty is then encoded by the variability of these cortical responses across time and across multiple trials, with higher uncertainty reflected in wider response distributions \cite{NIPS2002_a486cd07, delacour_neural_1997, orban_neural_2016}. Neural sampling thus provides an interpretation for the way the brain navigates uncertainty at the fundamental scale of neural computation.

Studies into neural sampling show that variability is not strictly a result of explicit stochastic noise, as traditionally believed. Instead, recent work demonstrates that strongly recurrent, balanced inhibitory-excitatory deterministic neural networks can generate the high-dimensional variability which stochastic models are built to capture \cite{terada_chaotic_2024, NIPS2002_a486cd07}. Terada and Toyoizumi link this sampling of a posterior distribution to Bayesian learning through training a recurrent network to sample static features and dynamic trajectories, generalising learned priors to novel, uncertain environments through biologically plausible learning \cite{terada_chaotic_2024}. Further biologically plausible implementations of such Bayesian sampling are energy-based models (EBMs) which suggest that fast-spiking interneurons can estimate the local partition function to enable fast, localised sampling dynamics \cite{dong_neural_nodate}. The sampling hypothesis has thus given a convincing theory of biological neural computation, inspired by the Bayesian brain framework (as described above in \ref{subsec:active_inf}), with testable and well-supported evidence that neural variability has an important functional role in computation \cite{orban_neural_2016, ujfalussy_sampling_2022}. Traditional Gaussian and Poisson noise models representing these stochastic network dynamics inherently correspond to continuous struggle to capture the full dynamics of complex interconnected inhibitory-excitatory (IE) biological neuronal network \cite{koren_efficient_2025, kilgore_biologically-informed_2025, song_training_2016}. Further study into the dynamics of these IE networks has thus motivated further functional models of neural activity.

\subsubsection*{Fractional neural sampling}

In reality, neural activity has been found to exhibit structure across multiple spatial and temporal scales, often described by scale-free statistics \cite{boonstra_scale-free_2013, he_scale-free_2014, he_temporal_2010, jiao_biologically_2026, jones_scale-free_2023, lin_scale-free_2016, mashour_conscious_2020, roberts_heavy_2015, zhang_heavy-tailed_2025}. This observation, together with the neural sampling hypothesis, motivates the design and theory of Fractional Neural Sampling (FNS) \cite{qi_fractional_2022}. FNS describes the complex spatiotemporal dynamics of spiking neural circuits, where localised activity patterns propagate through superdiffusive Lévy flights. Such Lévy flight dynamics of the centre of mass of the firings with an underdamped Langevin equation, which we shall call the FNS equations of motion,

\begin{equation}
\label{eq:fns}
\begin{aligned}
	d\textbf{x}_t &= \gamma b(\textbf{x}_t)dt + \beta \textbf{p}_t dt + \gamma^{1/\alpha} dL_t^{\alpha}, \\
	d\textbf{p}_t &= \beta b(\textbf{x}_t) dt
\end{aligned}
\end{equation}

where $\textbf{x}_t$ is the (usually 2 dimensional) vector of position, $\textbf{p}_t$ is a generalised auxiliary variable analagous to a momentum of the same dimensionality, $\beta > 0$ is a parameter defining the strength of coupling between $\textbf{p}_t$ and $\textbf{x}_t$, $\gamma > 0$ is the noise strength parameter, $ 0 < \alpha \leq 2$ is the Lévy tail exponent, and $dL_t^{\alpha}$ is the Lévy stable noise drawn from a $\mathcal{S}\alpha\mathcal{S}$ distribution. The term $b(\textbf{x}_t)$ captures a deterministic gradient term in the dynamics, defined such that the invariant distribution of $\textbf{x}_t$ after long time $t$ approaches the target distribution sampled, $\pi(\textbf{x})$. The exact expression for $b(\textbf{x})$ is derived in \cite{qi_fractional_2022}, and approaches the gradient of the logarithm of the density as $\alpha \rightarrow 2$, with the full form for any $ 0 < \alpha \leq 2$ given as

\begin{equation}
	b(\textbf{x}) = \frac{\mathcal{D}^{\alpha - 2}_{\textbf{x}}[\pi(\textbf{x}) \partial_{\textbf{x}} \pi(\textbf{x})]}{\pi(\textbf{x})}
\end{equation}

The FNS dynamics are called fractional, as the nonlocal dynamics of their distribution require fractional calculus to describe such spatial superdiffusion in Lévy flights \cite{tarasov_general_2021}. The theory of FNS thus unifies the scale-free nonlocal dynamics of neural firing patterns with the neural sampling framework derived from the ideal Bayesian scheme of active inference. The critical advantage of FNS over other standard continuous random walks is its ability to efficiently sample complex, multimodal probability distributions with far-apart modes—likely a major feature of processed natural environments replete with multiple distinct salient features. While Brownian-motion-based samplers typically remain trapped in local modes due to high potential energy barriers, the heavy-tailed, large jumps inherent in FNS enable the neural sampler to adaptively and freely switch between multiple modes \cite{neal_sampling_1996, latuszynski_mcmc_2025, simsekli_fractional_nodate, metzler_random_2000}. Furthermore, the generality of the FNS description permits the application of FNS equations of motion to gaze dynamics, as the motion of the eyes are driven by a IE neural circuit which is expected to keep key features of the FNS model invariant, such as superdiffusion, heavy-tailed statistics, and nonlocal interactions.

%This dynamical switching provides a functional explanation for the bimodality of estimate distributions observed in human perceptual switching and visual perception inference.
\subsection{Markov chain monte carlo methods}\label{subsec:MCMC}
This view of neural- and visual sampling of a target distribution motivates analogy to a foundational set of methods in machine learning. Markov Chain Monte Carlo (MCMC) methods have become the algorithmic basis for generating samples from target distributions in rigorous Bayesian inference in machine learning \cite{cranmer_advances_2023, latuszynski_mcmc_2025, speagle_conceptual_2020, andrieu_particle_2010}. MCMC operates by constructing a Markov chain, where the transition matrix is calculated from  whose equilibrium distribution is the desired posterior distribution. In the context of neural computation, biological circuits are often modelled as implementing continuous-time variants of MCMC, most notably Langevin sampling and Hamiltonian Monte Carlo (HMC) \cite{latuszynski_mcmc_2025, betancourt_conceptual_2018}. Langevin dynamics utilise local gradient information from the probability landscape, combined with Gaussian white noise, to explore the state space over time.  Despite their ubiquity, traditional MCMC methods encounter severe theoretical and practical limitations when exploring unknown, heavy-tailed, and/or multimodal target distributions. Standard Langevin and HMC dynamics are driven by continuous Brownian and Hamiltonian motion respectively, and the traversal of low-probability regions, equivalent to dynamical energy barriers, becomes exponentially slow . Consequently, the mean exit time for the sampler to leave one mode and discover another grows exponentially with the modal separation distance \cite{latuszynski_mcmc_2025}. To accelerate this sluggish exploration, algorithms often introduce auxiliary momentum variables. However, even with momentum-based acceleration, traversing deep energy valleys remains highly inefficient \cite{latuszynski_mcmc_2025, neal_sampling_1996, betancourt_conceptual_2018}. The fractional dynamics introduced in fractional sampling are one area of research into how fast mixing of samples in complex distributions can be achieved without the need for nontrivial parameter tuning. For instance, algorithms based on Hamiltonian Monte Carlo sampling and second-order Langevin sampling very similar to FNS dynamics have been briefly developed and explored \cite{simsekli_fractional_nodate, ijcai2018p419, qi_fractional_2022}. Under FNS dynamics, the mean exit time between modes scales linearly rather than exponentially with separation, demonstrating a profound computational advantage for efficient biological Bayesian inference \cite{qi_fractional_2022}.